\documentclass[11pt]{article}
\usepackage[a4paper,margin=2.5cm]{geometry}
\usepackage[T1]{fontenc}
\usepackage[utf8]{inputenc}
\usepackage{hyperref}

% Keep this value in sync with `project.version` in `pyproject.toml`.
\def\version{000.alpha.18}
\def\alphaht{\texttt{HydrologicalTwinAlphaSeries}}

\title{\alphaht\version User Guide\\Front API Only\\}
\author{hydrologicaltwin@minesparis.psl.eu}
\date{\today}

\begin{document}

\maketitle

\section{Purpose}

This document explains how to use HydrologicalTwinAlphaSeries (HTAS) from its
public front API only. The intended audience is an external consumer such as a
frontend, a workflow script, or a documentation example.

The only supported public entry point is the \texttt{HydrologicalTwin} facade
and its macro-methods:

\begin{itemize}
    \item \texttt{configure}
    \item \texttt{load}
    \item \texttt{describe}
    \item \texttt{fetch}
    \item \texttt{transform}
    \item \texttt{render}
    \item \texttt{export}
\end{itemize}

Do not drive HTAS through internal classes from \texttt{domain/},
\texttt{services/}, \texttt{tools/}, or private helper methods.

\section{Installation}

\begin{verbatim}
pixi install
\end{verbatim}

To activate the repository workflow described in the project README:

\begin{verbatim}
pixi run dev-setup
\end{verbatim}

Run the examples from the configured Pixi environment, for example through
\texttt{pixi run python} or from an activated Pixi shell.

\section{Public imports}

\begin{verbatim}
from HydrologicalTwinAlphaSeries import (
    ConfigGeometry,
    ConfigProject,
    HydrologicalTwin,
)
from HydrologicalTwinAlphaSeries.ht import (
    ConfigureRequest,
    DescribeRequest,
    FetchRequest,
    LoadRequest,
    RenderRequest,
    TransformRequest,
    ExportRequest,
)
\end{verbatim}

\section{Lifecycle}

The facade has an explicit lifecycle:

\begin{center}
\texttt{EMPTY} $\rightarrow$ \texttt{CONFIGURED} $\rightarrow$ \texttt{LOADED}
\end{center}

\begin{itemize}
    \item \texttt{configure} moves the twin from \texttt{EMPTY} to \texttt{CONFIGURED}.
    \item \texttt{load} moves the twin from \texttt{CONFIGURED} to \texttt{LOADED}.
    \item \texttt{describe}, \texttt{fetch}, \texttt{transform}, \texttt{render}, and
    \texttt{export} require at least the \texttt{LOADED} state.
\end{itemize}

If methods are called too early, HTAS raises \texttt{InvalidStateError}.

\section{Step 1 --- Build configuration objects}

HTAS expects a geometry configuration and a project configuration.

\begin{verbatim}
config_geom = ConfigGeometry.fromDict(
    {
        "ids_compartment": [1],
        "resolutionNames": {1: [["AQ_LAYER"]]},
        "ids_col_cell": {1: 0},
        "obsNames": {},
        "obsIdsColCells": {},
        "obsIdsColNames": {},
        "obsIdsColLayers": {},
        "obsIdsCell": {},
        "extNames": {},
        "extIdsColNames": {},
        "extIdsColLayers": {},
        "extIdsColCells": {},
    }
)

config_proj = ConfigProject.fromDict(
    {
        "json_path_geometries": "geometry.json",
        "projectName": "demo",
        "cawOutDirectory": "/path/to/CAWAQS/OUTPUTS",
        "startSim": 2000,
        "endSim": 2001,
        "obsDirectory": "/path/to/OBS",
        "regime": "annual",
    }
)
\end{verbatim}

\section{Step 2 --- Create and configure the facade}

You can either configure the twin at construction time or call
\texttt{configure()} explicitly.

\subsection*{Constructor path}

\begin{verbatim}
twin = HydrologicalTwin(
    config_geom=config_geom,
    config_proj=config_proj,
    out_caw_directory=config_proj.cawOutDirectory,
    obs_directory=config_proj.obsDirectory,
    temp_directory="/path/to/TEMP",
    metadata={"consumer": "front-api-example"},
)
\end{verbatim}

\subsection*{Explicit request path}

\begin{verbatim}
twin = HydrologicalTwin()
twin.configure(
    request=ConfigureRequest(
        config_geom=config_geom,
        config_proj=config_proj,
        out_caw_directory=config_proj.cawOutDirectory,
        obs_directory=config_proj.obsDirectory,
        temp_directory="/path/to/TEMP",
        metadata={"consumer": "front-api-example"},
    )
)
\end{verbatim}

\section{Step 3 --- Load the project}

The production path is to load the configured compartments through the public
\texttt{LoadRequest}. A frontend should stay at this facade level and must not
orchestrate internal backend objects directly.

\begin{verbatim}
twin.load(
    request=LoadRequest(
        ids_compartments=[1],
        geo_provider=geo_provider,
    )
)
\end{verbatim}

For smoke tests that validate only the lifecycle, an empty load is acceptable:

\begin{verbatim}
twin.load(request=LoadRequest(compartments={}))
\end{verbatim}

\section{Step 4 --- Inspect the available catalog}

Use \texttt{describe()} immediately after loading to discover what the frontend
can ask for next.

\begin{verbatim}
description = twin.describe(request=DescribeRequest())

print(description.state)
print(description.n_compartments)
print(description.catalog.extract_kinds)
print(description.catalog.transform_kinds)
print(description.catalog.render_kinds)
\end{verbatim}

\section{Step 5 --- Fetch data through stable workflow kinds}

\subsection*{Simulation matrix}

\begin{verbatim}
values = twin.fetch(
    request=FetchRequest(
        kind="simulation_matrix",
        id_compartment=1,
        outtype="MB",
        param="recharge",
        syear=2000,
        eyear=2001,
        id_layer=0,
    )
)

print(values.data.shape)
print(values.dates[:3])
\end{verbatim}

\subsection*{Observations}

\begin{verbatim}
observations = twin.fetch(
    request=FetchRequest(
        kind="observations",
        id_compartment=1,
        syear=2000,
        eyear=2001,
    )
)
\end{verbatim}

\subsection*{Frontend bundle}

\begin{verbatim}
bundle = twin.fetch(
    request=FetchRequest(
        kind="sim_obs_bundle",
        id_compartment=1,
        outtype="Q",
        param="discharge",
        syear=2000,
        eyear=2001,
        compute_criteria=True,
    )
)
\end{verbatim}

Other public fetch kinds are exposed by \texttt{description.catalog.extract\_kinds}.

\section{Step 6 --- Transform fetched payloads}

\subsection*{Temporal aggregation}

\begin{verbatim}
monthly = twin.transform(
    request=TransformRequest(
        kind="temporal_aggregation",
        data=values.data,
        dates=values.dates,
        frequency="Monthly",
        agg_dimension="mean",
    )
)
\end{verbatim}

\subsection*{Performance criteria}

\begin{verbatim}
criteria = twin.transform(
    request=TransformRequest(
        kind="criteria",
        bundle=bundle,
        metrics=["nash", "bias", "rmse"],
    )
)
\end{verbatim}

\subsection*{Budget and hydrological regime}

\begin{verbatim}
budget = twin.transform(
    request=TransformRequest(
        kind="budget",
        id_compartment=1,
        param="recharge",
        sdate=2000,
        edate=2001,
        frequency="Annual",
        agg_dimension="mean",
    )
)

regime = twin.transform(
    request=TransformRequest(
        kind="hydrological_regime",
        id_compartment=1,
        outtype="Q",
        param="discharge",
        sdate=2000,
        edate=2001,
    )
)
\end{verbatim}

\section{Step 7 --- Render artefacts}

\begin{verbatim}
budget_plot = twin.render(
    request=RenderRequest(
        kind="budget_barplot",
        data={
            "data": budget.data,
            "date_labels": budget.date_labels,
            "param": budget.param,
        },
        plot_title="Recharge budget",
        output_folder="/path/to/output",
        output_name="recharge_budget",
        yaxis_unit="mm/j",
    )
)

regime_plot = twin.render(
    request=RenderRequest(
        kind="hydrological_regime",
        data=regime.data,
        obs_point_names=regime.obs_point_names,
        month_labels=regime.month_labels,
        var="discharge",
        units="m3/s",
        savepath="/path/to/output/regime",
        interactive=False,
        staticpng=True,
        staticpdf=True,
    )
)

print(budget_plot.artefacts)
print(regime_plot.artefacts)
\end{verbatim}

\section{Step 8 --- Export the twin}

\begin{verbatim}
export_result = twin.export(
    request=ExportRequest(
        path="/path/to/output/twin.pkl",
        fmt="pickle",
    )
)

print(export_result.path)
\end{verbatim}

\section{Recommended frontend workflow}

\begin{enumerate}
    \item Build \texttt{ConfigGeometry} and \texttt{ConfigProject}.
    \item Create one \texttt{HydrologicalTwin} instance per project/session.
    \item Call \texttt{configure}.
    \item Call \texttt{load}.
    \item Call \texttt{describe} to populate UI selectors and supported workflow kinds.
    \item Use \texttt{fetch} for raw frontend-ready payloads.
    \item Use \texttt{transform} for computations derived from fetched data.
    \item Use \texttt{render} only when file artefacts are explicitly needed.
    \item Use \texttt{export} to persist the loaded twin snapshot.
\end{enumerate}

\section{What the frontend should not do}

\begin{itemize}
    \item Do not import from \texttt{domain/}, \texttt{services/}, or private helpers.
    \item Do not manipulate \texttt{Compartment}, \texttt{Mesh}, \texttt{Observation}, or extraction objects directly.
    \item Do not skip lifecycle order.
    \item Do not rely on undocumented internal methods such as \texttt{read\_values()}.
\end{itemize}

\section{Minimal smoke test}

\begin{verbatim}
twin = HydrologicalTwin()
twin.configure(
    request=ConfigureRequest(
        config_geom=config_geom,
        config_proj=config_proj,
        out_caw_directory=config_proj.cawOutDirectory,
        obs_directory=config_proj.obsDirectory,
    )
)
twin.load(request=LoadRequest(compartments={}))
description = twin.describe()
assert description.state == "LOADED"
\end{verbatim}

\end{document}
